% ============================================================
\section{Introdução e objetivos}\label{sec:introducao}
% ============================================================

Um modelo de linguagem do tipo Transformer lê o texto como uma sequência de \emph{tokens}
(palavras ou pedaços de palavras) e representa cada um por um vetor. Na entrada, esse vetor é
uma linha de uma tabela fixa, a matriz de \emph{embeddings}: todas as ocorrências de
\emph{banco} começam com o mesmo vetor, seja num texto sobre finanças, sobre computação ou
sobre geografia. Ao atravessar os blocos do modelo, cada vetor recebe informação dos tokens
anteriores, e as ocorrências de um mesmo token deixam de ser iguais; no último bloco, o vetor
serve para prever o \emph{próximo} token. Estudos anteriores mediram essa passagem com médias
sobre pares de vetores, como o cosseno médio entre as ocorrências de uma
palavra~\cite{ethayarajh2019}. Essas médias dizem quanto os vetores se afastam, mas não dizem
\emph{quem é vizinho de quem}, nem se a mudança é só local (troca de alguns vizinhos) ou se
altera a organização em escala maior (\emph{hubs}, distâncias, comunidades).

Este projeto descreve essa passagem com redes complexas. Cada vértice é uma \emph{ocorrência}
de um token num parágrafo, e cada vértice aponta para os $k=10$ vértices de vetor mais parecido
(uma rede de $k$ vizinhos mais próximos, k-NN). Os mesmos vértices recebem quatro vetores do
modelo Qwen3-4B-Base: o \emph{embedding} de entrada (rede \cod{lex}) e as saídas dos blocos 1,
18 e 36 (redes \cod{L01}, \cod{L18} e \cod{L36}). Como as redes têm os mesmos vértices e
diferem só nos arcos, é possível comparar a vizinhança de cada ocorrência entre as camadas,
separar o que muda em volta de cada vértice do que muda na rede inteira e comparar as
comunidades com rótulos conhecidos dos vértices (o token, o tema, o artigo, o parágrafo, o
próximo token). Isso transforma ``a representação ficou contextual'' numa afirmação mensurável.

\subsection{Objetivos}\label{sec:objetivos-def}

\textbf{Objetivo geral} (da proposta): caracterizar, por meio de redes complexas, quanto e de
que forma a estrutura de similaridade entre representações de tokens se transforma entre o
espaço lexical de entrada e as representações contextuais produzidas pelas camadas inicial,
intermediária e final de um Transformer.

\textbf{Objetivos específicos}: responder a quatro perguntas, cada uma com uma hipótese
formulada antes dos resultados.
\begin{description}
    \item[P1. Quanto e de que forma a organização muda?] Medir a troca de vizinhos de cada
    ocorrência, a fração de vizinhos do mesmo token (dominância lexical) e as métricas globais
    (graus de entrada, clusterização, distâncias, componentes, comunidades). \emph{H1:} a
    dominância lexical cai, mas continua acima do acaso no bloco 36; surgem \emph{hubs}; as
    comunidades se alinham menos ao token.
    \item[P2. Em que camadas a mudança se concentra?] Comparar as transições entre as redes e
    acompanhar as medidas nos 36 blocos, por faixa de frequência e categoria de token.
    \emph{H2:} pedaços de palavra, que precisam dos vizinhos para formar uma palavra, mudam
    antes das palavras inteiras.
    \item[P3. Surgem comunidades associadas a contextos semelhantes?] Comparar as comunidades
    com o token, o tema, o artigo, o parágrafo, a sentença e o próximo token. \emph{H3:} a
    concordância com o token cai e a com o contexto sobe; no último bloco, cresce a com o
    próximo token.
    \item[P4. As ocorrências de um mesmo token se separam conforme o contexto?] Medir se as
    ocorrências de palavras com sentidos diferentes em temas diferentes (\emph{banco},
    \emph{órgão}, \emph{matriz}\ldots), idênticas na entrada, passam a ocupar regiões e
    comunidades diferentes. \emph{H4:} elas se separam, por tema e por sentido, mais do que
    palavras de sentido estável (\emph{ano}, \emph{grande}, \emph{água}\ldots).
\end{description}
A proposta listava também quatro etapas: construir as redes com cerca de $1{,}5\times10^4$
ocorrências e os mesmos vértices em todas; calcular as métricas estruturais pedidas na
disciplina; reportar os resultados por faixa de frequência e categoria de token, com testes de
robustez; e visualizar as redes.

\textbf{Organização.} A Seção~\ref{sec:parcial} resume o projeto e a entrega parcial; as
Seções~\ref{sec:metodologia} e~\ref{sec:resultados} trazem a metodologia e os resultados; a
Seção~\ref{sec:objetivos} mostra como cada objetivo foi cumprido; e a
Seção~\ref{sec:discussao} discute os resultados e as limitações. O código da implementação está
em \url{\repo}.
